\medskip

Dans un repère orthonormé \Oijk{} on considère les points :
\begin{center} A$(1~;~1~;~-4)$,\quad B$(2~;~-1~;~-3)$,\quad  C$(0~;~-1~;~-1)$\quad  et\quad  $\Omega$(1~;~1~;~2).\end{center}

\smallskip

\begin{enumerate}
\item Démontrer que les points A, B et C définissent un plan.
\item 
	\begin{enumerate}
		\item Démontrer que le vecteur $\vect{n}$ de coordonnées $\begin{pmatrix}1\\1\\1\end{pmatrix}$ est normal au plan (ABC).
		\item Justifier qu'une équation cartésienne du plan (ABC) est $x + y + z + 2 = 0$.
	\end{enumerate}
\item 
	\begin{enumerate}
		\item Justifier que le point $\Omega$ n'appartient pas au plan (ABC).
		\item Déterminer les coordonnées du point H, projeté orthogonal du point $\Omega$ sur le plan (ABC).
	\end{enumerate}
			\end{enumerate}

On admet que $\Omega \text{H} = 2\sqrt 3$.

On définit la sphère $S$ de centre $\Omega$ et de rayon $2\sqrt 3$ comme l'ensemble de tous les points M de l'espace tels que $\Omega\text{M} = 2\sqrt 3$.

	\begin{enumerate}[resume]

\item Justifier, sans calcul, que tout point N du plan (ABC), distinct de H, n'appartient pas à la sphère $S$.

On dit qu'un plan $\mathcal{P}$ est tangent à la sphère $S$ en un point K lorsque les deux conditions suivantes sont vérifiées :

\[\begin{array}{l}
\bullet~~ \text{K} \in \mathcal{P} \cap S\\
\bullet~~ (\Omega\text{K}) \perp \mathcal{P}
\end{array}
\]
\item Soit le plan $\mathcal{P}$ d'équation cartésienne $x +y - z - 6 = 0$ et le point K de coordonnées K(3~;~3~;~0).

Démontrer que le plan $\mathcal{P}$ est tangent à la sphère $S$ au point K.
\item On admet que les plans (ABC) et $\mathcal{P}$ sont sécants selon une droite $(\Delta)$.

Déterminer une équation paramétrique de la droite $(\Delta)$.
\end{enumerate}

\bigskip
